\section{Failure Attribution}

Canvass arithmetic is useful only if failures are attributed to the right layer.
RCOUNT transcripts therefore name equation ids rather than producing only a
single pass/fail bit.

\begin{center}
\small
\begin{tabularx}{\textwidth}{lXX}
\toprule
Layer & Failure & Interpretation \\
\midrule
Contest sum & Candidate votes plus residuals do not equal counted ballots & Local summary arithmetic is wrong. \\
Jurisdiction total & Jurisdiction total does not equal reporting-unit summaries & Rollup is wrong for that contest/status. \\
Canvass correction & Changed canvass lacks public correction event & Transition evidence is missing. \\
Batch summary & Batch summary references absent manifest & Accounting evidence is missing. \\
Accepted ballots & Accepted ballots differ from counted plus rejected & Batch manifest accounting is inconsistent. \\
Source hash & Raw source file changes after packaging & Source evidence no longer matches the package index. \\
\bottomrule
\end{tabularx}
\end{center}

In verifier transcripts these layers correspond to stable machine ids:
\begin{quote}
\small
\begin{verbatim}
contest_selection_sum
jurisdiction_contest_total
canvass_correction_event
batch_summary_total
accepted_ballots
source_hash_match
\end{verbatim}
\end{quote}

This attribution matters rhetorically as well as technically. A source hash
failure is not the same as a bad jurisdiction total. A missing batch manifest is
not the same as a fraudulent canvass. RCOUNT's contribution is to make those
distinctions visible enough that later legal, administrative, or audit decisions
can reason from a stable mechanical record.
